\begin{flushleft}
\textbf{Interview mit einem Stadtplaner über digitale Beteiligung und Regeln für öffentliche Räume} \\
Durchgeführt vom Autor im Rahmen der Masterarbeit. Das Interview wurde transkribiert und sprachlich leicht geglättet, ohne inhaltliche Änderungen.
\end{flushleft}

\begin{description}
\item[Author:] Welche Form von partizipativer Stadtplanung begegnen Ihnen am meisten? Oder was für digitale Werkzeuge werden genutzt?
\item[Stadtplaner:] Es gibt ja unterschiedliche Arten von Partizipation. Das eine ist Information, dann eben Mitgestaltung oder auch Mitentscheidung. Und es gibt, je nachdem, in was für einem Kontext man sich befindet, eine formelle Art der Beteiligung. Das heißt, es gibt eine Information darüber, wenn ein Bebauungsplan erstellt wird. Und das wird dann, in Berlin beispielsweise, auf einer Plattform angekündigt, und man kann dann Kommentare dazu schreiben. Das ist eine Form, die relativ formalisiert ist. Das heißt, dass die bei sehr vielen Beteiligungsverfahren, die das Land Berlin durchführt, auch angewandt wird. Das hat so ein paar Fallstricke, aber ist zumindest immer eine Möglichkeit, digitale Beteiligung durchzuführen. Alles andere sind dann eher experimentelle Verfahren, die mal ausprobiert werden. Also dass man irgendwelche, weiß ich nicht, virtuellen Szenarien entwirft mit den Bürgerinnen und Bürgern. Das ist aber nicht die Standardbeteiligung. Also wie gesagt, die Standardbeteiligung, die die Stadtplanungsämter durchführen, ist plattformbasiert. Dann wird ein Plan ausgelegt, kennen Sie sicher. Und dann gibt es die Möglichkeit, da Kommentare zu geben.
\item[Author:] Haben Sie da irgendwelche Erfahrungen, was denn am meisten als Kommentar gesetzt wird? Gibt es da irgendwelche Muster?
\item[Stadtplaner:] Nein, das hängt natürlich vom Verfahren ab, um was es geht. Also es gibt da ganz unterschiedliche Anlässe, um das zu machen. Und je nachdem, was für ein Anlass es ist, sind natürlich auch die Kommentare unterschiedlich. Wenn es darum geht, dass ein neuer Bebauungsplan aufgestellt wird, dann ist eines der häufigsten Kommentare, auch wenn ich das jetzt nicht statistisch validiert habe, dass es die Bürgerinnen und Bürger generell nicht gut finden und erstmal grundsätzlich ablehnen, weil das führt zu mehr Verkehr, das führt zu Staub und Lärm beim Bauprozess, und eventuell wird auch noch die Sicht eingeschränkt. Also oftmals so eine ablehnende Haltung. Natürlich wird das dann auch teilweise noch differenziert. Aber das ist oftmals ein Ergebnis von Beteiligungsprozessen, dass das Projekt eben grundsätzlich infrage gestellt wird.
\item[Author:] Um mal so ein bisschen mehr zu diesem vielleicht auch experimentellen Teil zu kommen. Ich glaube, Sie hatten da auch bei ein paar solcher experimentellen Szenarien mitgewirkt, wenn ich das richtig im Kopf habe?
\item[Stadtplaner:] Ja, wir hatten auch ein Forschungsprojekt, wo es um Räume ging, die geteilt werden sollen. Und da hatten wir auch so ein paar Sachen ausprobiert, also auch so ein Beteiligungsformat. Das war ein bisschen schwierig, weil das war dann zu Corona Zeiten und insofern war es dann eben nur im digitalen Raum. Wir haben dann auch so ein paar Nachbarschaftstreffen gemacht. Da ging es darum, dass wir erstmal wissen wollten von den Bewohnerinnen und Bewohnern, welche Orte im Stadtteil ihnen besonders wichtig sind, welche Wünsche sie haben, was für Ängste sie auch haben und wo man vielleicht mit Flächen neu umgehen kann. Wir haben das aber relativ, ich sage mal, traditionell dann gemacht in Form wie gesagt einer Online-Befragung. Und daraufhin hatten wir dann ein beteiligtes Architekturbüro, hat dann daraufhin Entwürfe gezeichnet, und diese Entwürfe haben wir dann noch mal mit den Bewohnergruppen abgestimmt. Also es war nicht so, dass wir – was Ihr Forschungsprojekt macht – die Entwürfe im digitalen Raum irgendwie zur Verfügung gestellt haben, sondern es waren eben Zeichnungen, und über die haben wir dann diskutiert. Wir hatten noch ein anderes ganz nettes Format, wo wir einen Zeichner hatten, der Live-Zeichnungen gemacht hat. Das heißt, es gab ein Fest, die Leute haben erzählt, was sie sich alles wünschen, und dann gab es einen, der das live gezeichnet hat. Das war für die Bewohner eigentlich ganz nett anzugucken - okay, wie sieht sowas denn dann aus? Das waren dann auch teilweise ein bisschen verrückte Sachen. Da gab es dann eine Gruppe, die wollte einen öffentlichen Pizzabackofen, hat er das gezeichnet, und dann konnte man sich anschauen, wie so eine Skizze dann aussieht. Es war aber nicht so, dass wir von der Stadtverwaltung beauftragt wurden, Beteiligungsverfahren durchzuführen.
\item[Author:] Wenn wir jetzt daran denken, dass Bürger über so ein System frei vorschlagen können. Gibt es da bestimmte Aspekte, die Ihnen einfallen würden, wo so etwas am ehesten scheitern könnte oder falsch sein könnte? So etwas wie Abstände, Barrierefreiheit etc.?
\item[Stadtplaner:] Ja, der Hauptpunkt ist Eigentumsfragen [EL-IN-001]. Wem gehört der Grund und Boden? Oftmals hat man da Vorstellungen darüber, dass man irgendwie etwas machen will, aber das ist dann auf dem Grundstück von jemand anderem, und dann funktioniert es eben nicht. Also das heißt, diese Eigentumsfragen spielen eine ganz große Rolle. Wenn man sich jetzt auf öffentlichen Plätzen bewegt, dann ist es einfacher. Das ist oftmals das Thema, dass diese Eigentumsfragen vernachlässigt werden, aber auch Haftungsfragen. Wenn man jetzt irgendwie den Schulhof neu umgestalten möchte, dann ist das ein Thema. Wie gesagt, hängt mit Eigentum auch zusammen. Ansonsten sind so ein paar Punkte. Klar, Abstandsflächen [EL-IN-002]. Wenn ich jetzt sage, ich will ein neues Stadion – also wenn ich jetzt im großen Maßstab denke, ich will irgendwie ein neues Stadion bei mir im Kiez bauen – dann ist es aus 1000 Gründen nicht möglich. Also die Maßstäblichkeit, sich das zu überlegen [EL-IN-003]. Technische Umsetzung: Wenn ich Kinder frage, was wollen sie am liebsten? Dann sagen die, sie wollen irgendwie ein Schwimmbad auf der Straße. Das ist dann natürlich technisch nicht umsetzbar. Ansonsten könnten Hindernisse auch noch irgendwie Lärmbelästigung sein. Wenn ich daran denke, dass Spielplätze, Sportplätze und so weiter gewünscht werden, dass das dann zu einer Lärmbelästigung führt und daher nicht möglich ist. Dann noch ganz wichtig ist die Frage der Pflege. Also wenn ich jetzt sage, ich möchte hier eine Grünfläche haben, dann stellt sich oftmals die Frage, wer pflegt denn, wer ist für den Unterhalt zuständig, wer entfernt Müll usw.? Und das wird oftmals auch vernachlässigt oder könnte auch ein Hindernis sein, gerade bei öffentlichen Beteiligungsverfahren. Weil ich natürlich sagen kann, wir haben kein Geld, um noch mehr Grünflächen zu pflegen.
\item[Author:] Und bei Eigentumsfragen, wem gehört was? Schaut man dann als Stadtplaner einfach ins Grundbuch?
\item[Stadtplaner:] Das ist nicht so einfach. Die öffentliche Hand hat natürlich Zugriff auf Grundbücher und weiß, wem was gehört. Diese Daten sind aber nicht öffentlich zugänglich. Das heißt, wenn ich Leute beteiligen möchte, dann kann ich nur sagen, das ist jetzt Privatgrundstück, das heißt, hier kann nicht beteiligt werden, und hier kann jetzt irgendwie der Platz umgestaltet werden. Es sei denn, ich habe vorher irgendwie einen Deal mit dem Eigentümer, dass der sagt, ja, ich habe auch Interesse daran, dass der Platz umgestaltet wird. Anders gesagt, ich kann beim Beteiligungsverfahren sagen, was ist öffentlich und was ist privat. Ich darf aber nicht automatisch auch sagen, wem die private Fläche dann gehört.
\item[Author:] Gibt es da irgendwelche Standards, irgendwelche Checklisten, auf die man sich als Stadtplaner stützt?
\item[Stadtplaner:] Es gibt die Berliner Leitlinien für Partizipation. Das wäre so etwas, wo man sich daran stützen kann. Ansonsten, was immer ein Thema ist, ist die Frage des Zwecks der Beteiligung. Also worum geht es eigentlich und was passiert mit den Beteiligungsergebnissen? Geht es nur darum, Ideen zu sammeln? Oder geht es darum, wirklich etwas umzusetzen? Die Gefahr bei Beteiligungsprozessen ist eigentlich immer die, dass die Bürgerinnen und Bürger dann hinterher das Gefühl haben, dass nichts passiert, obwohl sie so viele tolle Ideen gemacht haben. Also das heißt, dieses Nachverfolgen, was passiert mit meinen Ideen, warum wird vielleicht auch was nicht umgesetzt, das ist etwas, was wichtig ist und was oftmals vernachlässigt wird. Weil, wenn ich mich jetzt als Bürger hinsetze und sage, ich möchte etwas, das heißt, ich investiere da auch Zeit, dann möchte ich natürlich wissen, was passiert mit meinen Ergebnissen.
\item[Author:] Wenn man jetzt an so ein Tool oder Mein Berlin denkt, was sind da so Voraussetzungen, damit eine Behörde so ein Tool überhaupt wirklich ernst nehmen kann?
\item[Stadtplaner:] Natürlich sind bei Planungsprozessen die Bürger nur ein Akteur, der zu beteiligen ist. Man hat dann, die Träger öffentlicher Belange, das heißt, Feuerwehr muss beteiligt werden, es müssen verschiedene Ämter beteiligt werden. Es muss, wie gesagt, wenn man unterschiedliche Eigentümer hat, die noch beteiligt werden. Das heißt, es ist selten der Fall, dass die Einwände der Bürger etwas ganz Neues bringen, was im Beteiligungsprozess nicht schon von anderen angesprochen wurde. Also, insofern glaube ich, dass es sicherlich hilfreich ist für die Stadtverwaltung, wenn man da mit einem neuen Tool noch besser die Meinungen der Bürger abholen könnte, aber ob das jetzt wirklich dazu führt, dass die Beteiligungsprozesse besser werden? Das könnte ich mir vorstellen, dass es ohnehin so komplex ist und dass jede Lösung nur ein Kompromiss ist. Vielleicht wäre das interessant für ein Planungstool, dann zu sagen, aus welchen Gründen haben wir uns für eine Lösung entschieden [EL-IN-004]. Also diese Nachvollziehbarkeit, das könnte ich mir vorstellen, dass das ein Mehrwert ist.
\item[Author:] Wenn man sich jetzt vorstellt, dass so sehr kleine Einwendungen, etwa eine neue Bank, ist. Denken Sie, dass das irgendwie schon noch realistisch ist, dass da die Behörde dann sagen könnte, dann bauen wir die Bank da noch mehr hin?
\item[Stadtplaner:] Ja, das könnte ich mir schon vorstellen. Also, vielleicht bei konkreten Gestaltungsfragen, wo es jetzt nicht darum geht, ob in meinem Hinterhof jetzt ein fünfgeschossiges Haus gebaut wird, sondern wenn es um so Kleinigkeiten geht, wo können wir eine neue Bank hinstellen? Wo brauchen wir mehr Schattenplätze oder irgendwie so was? Da könnte ich mir schon vorstellen, dass die Einwände der Bürger dann auch leichter umsetzbar sind. Aber da müsste es dann so gestaltet sein, dass das Verfahren relativ unaufwendig für die Verwaltung ist. Also wenn ich mir überlege, wie überlastet die Verwaltung ist, und dann zu sagen, jetzt machen wir ein Beteiligungsverfahren, wo es nur darum geht, ob wir eine oder drei Bänke hinstellen sollen, dann wird die Verwaltung vielleicht sagen, das ist jetzt nicht so ganz entscheidend, wir haben andere Probleme, die viel größer sind. Also, daher die Möglichkeit, so was einfach zu realisieren, ist für die Verwaltung etwas, was das Programm stärken würde. Es gibt mit Mein Berlin eine Plattform, über die die Beteiligung abläuft. Und wenn es jetzt darum geht, Beteiligungsverfahren oder neue Tools zu entwickeln, dann denke ich, dass diese Einbindung in Mein Berlin ein wichtiger Faktor ist, weil dadurch es für die Verwaltung dann auch einfacher ist, so was umzusetzen [EL-IN-005]. Und auf der anderen Seite durch die Plattform eine gewisse Öffentlichkeit sowieso entsteht, weil auf Mein Berlin kann man alles Mögliche suchen. In anderen Ländern, in anderen Städten mag das anders sein, aber für Berlin glaube ich, dass eben diese Verknüpfung mit der existierenden Plattform wichtig ist.
\item[Author:] Sie irgendwelche Risiken oder Probleme, wenn man so statische, feste Regeln in so ein System überführt?
\item[Stadtplaner:] Ja, also das Risiko ist aus Bürgersicht, dass dann gesagt wird, es interessiert gar nicht die Meinung der Bürger, ob was gut ist oder was schlecht ist, sondern es geht ja nur darum, die Farbe der Mauer oder die Farbe des neuen Gebäudes zu bestimmen. Also, je gestrickter es ist, umso weniger Entscheidungsspielraum gibt es natürlich. Und teilweise besteht dann die Gefahr, dass das als Alibi-Beteiligung wahrgenommen wird. Also, wenn ich sage, ich möchte jetzt hier eine neue Straße bauen und die Bürger können entscheiden, ob der Beton hellgrau oder dunkelgrau ist. Das glaube ich, ist eine Gefahr. Auf der anderen Seite, wäre es natürlich durchaus hilfreich, wenn man die Kosten irgendwie im Blick hat. Wenn ich sage, ich möchte jetzt etwas haben, aber ich habe ein Budget von zwanzigtausend Euro, was ist denn damit möglich? Und vielleicht wären da digitale Tools irgendwie ein Ansatz zu sagen, ich kann alles machen, aber es darf nicht mehr als zwanzigtausend Euro kosten. Also es ist immer so ein bisschen, es ist immer eine Balance zwischen, ich will nicht, dass da nur Sachen rauskommen, die komplett utopisch sind und sowieso nicht realisiert werden können, auf der einen Seite. Und auf der anderen Seite eben so eine Alibi-Beteiligung, wo es eigentlich um gar nichts mehr geht und alles schon entschieden ist. Also, das ist immer eine Abwägung.
\item[Author:] Denken Sie, dass noch irgendwelche offenen Fragen sind, die in so einem Projekt unbedingt berücksichtigt werden sollten?
\item[Stadtplaner:] Ich glaube, der Punkt ist bei der Bürgerbeteiligung, dass die Herausforderungen, die existieren, gerade im städtischen Raum - der Raum ist knapp und der Raum ist umkämpft. Natürlich sagt die Stadtplanung, am liebsten wäre uns, wenn wir hier toll Bänke für die Bürger aufstellen können. Auf der anderen Seite ist das aber genau die Fläche, die die Feuerwehr braucht, um den Wendekreis von ihren Löschfahrzeugen zu benutzen. Und gleichzeitig ist das eine Fläche, die mit Altlasten verseucht ist. Das heißt, wo sich die Bürger eigentlich ohnehin nicht aufhalten sollten. Diese Komplexität, die halt dazu führt, dass selten genau das umgesetzt wird, was die Bürgerinnen und Bürger sich wünschen. Das ist, glaube ich, die Herausforderung. Aber, wenn man das als Potenzial sehen will, dann ist das ja auch durch die digitalen Tools vielleicht möglich, diese Komplexität auch dann besser aufzuzeigen und die unterschiedlichen Anforderungen an den Raum auch zu beleuchten [EL-IN-004]. Wenn ich sage, ich wünsche mir jetzt hier für meine Straße eine Bank mit Schatten und einem Trinkbrunnen nebendran, weil alle meine Nachbarn sagen auch, dass das total notwendig ist und wir schlagen das gemeinsam vor und es kommt dann eben nicht, dann muss ich als Bürger abgeholt werden und mir muss erklärt werden, warum das nicht passiert. Und das ist so ein bisschen die Krux bei der Beteiligung. Vielleicht kann man es sicherlich noch nicht so eins zu eins komplett umsetzen, aber zumindest ein paar. Und das versuchen Sie ja auch mit Ihren Regeln, die Sie da aufstellen wollen. Also wie gesagt, ich glaube, es geht nicht nur darum, den Bürgern irgendwie ein Tool zu geben, wo diese alles Mögliche einzeichnen können, sondern den Bürgern die Möglichkeit zu geben, auch zu sagen, warum eben die Sachen vielleicht dann doch nicht so einfach realisierbar sind. Und wenn es eine Geldsache ist oder irgendwie eben die Grenzen der Beteiligung aufzeigen. Sie werden immer noch über die Verwaltung schimpfen, aber vielleicht werden sie sehen, es ging jetzt halt nicht anders aus den und den Gründen.
\end{description}